Die Arbeit ist wie folgt aufgebaut:

Kapitel \ref{ch:grundlagen} vermittelt die theoretischen Grundlagen zu Leichter Sprache, Sprachmodellen und Methoden der Textbewertung. 
Kapitel \ref{ch:analysephase} analysiert die bestehende SWR-Pipeline sowie die funktionalen und nichtfunktionalen Anforderungen an das geplante System.
Kapitel \ref{ch:konzeption} formuliert das konzeptionelle Design des Evaluations- und Verbesserungssystems, während die konkrete Implementierung in Kapitel \ref{ch:Implementierung} kurz erläutert wird. 
Kapitel \ref{ch:evaluation} widmet sich der Evaluation der Ergebnisse anhand von Testfällen und Vorher-Nachher-Vergleichen. 
Abschließend werden die Ergebnisse in Kapitel \ref{ch:diskussion} diskutiert und in Kapitel \ref{ch:fazit} zusammengefasst sowie ein Ausblick auf zukünftige Weiterentwicklungen gegeben.